\chapter{The Standard Library, Part 1}
\label{ch:stdlib-1}

A language is only as productive as the library that comes with it. Salam ships a
substantial standard library so that you rarely have to reinvent the basics. This
chapter tours the everyday modules input and output, formatting, mathematics,
randomness, and time with runnable examples of the functions you will use most.
The next chapter covers data structures and the world outside your program.

\section{Input and output: \texttt{io}}

You already know \code{println} and \code{print}, which are built in. The
\code{io} module offers the same operations as named functions (handy when you
import it for file work) plus the all-important \code{Input}, which reads a line
of text the user types:

\begin{salam}
import "io"

func main:
    io.Print("What is your name? ")
    name str = io.Input()
    io.Println("Hello, " + name + "!")
end
\end{salam}

\code{io.Input()} reads one line from the keyboard and returns it as a \code{str}
(without the trailing newline). Combined with the \code{str} conversions from
Chapter~8, this is how you build interactive programs:

\begin{salam}
import (
    "io"
    "str"
)

func main:
    io.Print("Enter a number: ")
    n auto = str.ToInt(io.Input())
    io.Println("Its square is " + str.FromInt((n * n) as i64))
end
\end{salam}

\begin{note}
Because \code{io.Input} returns text, you must convert it to a number with
\code{str.ToInt} (or \code{str.ToFloat}) before doing arithmetic. Reading input as
text and converting explicitly is safer than guessing types you decide exactly
how each line should be interpreted.
\end{note}

\section{Formatting values: \texttt{fmt}}

The \code{fmt} module turns values into strings and lays them out neatly. The
conversions are the bread and butter:

\begin{salam}
import "fmt"

func main:
    println fmt.Int(255 as i64)        // "255"
    println fmt.Float(3.14159)         // "3.14159"
    println fmt.Bool(true)             // "true"
end
\end{salam}

\begin{progout}
255
3.14159
true
\end{progout}

For aligning text into columns tables, menus, reports \code{fmt} provides
padding and repetition:

\begin{salam}
import "fmt"

func main:
    println fmt.PadLeft("hi", 5)       // "   hi"
    println fmt.PadRight("hi", 5) + "|" // "hi   |"
    println fmt.Repeat("=", 10)        // "=========="
end
\end{salam}

\begin{progout}
   hi
hi   |
==========
\end{progout}

\code{PadLeft} and \code{PadRight} widen a string to a given width by adding
spaces on the left or right, so numbers and labels line up in columns.
\code{Repeat} builds a string of repeated copies perfect for rules and
separators.

\section{Mathematics: \texttt{math}}

The \code{math} module covers the numeric functions you would expect, working on
\code{f64} values (with a few integer helpers). Here are the most useful, grouped
by purpose.

\subsection{Rounding and sign}

\begin{salam}
import "math"

func main:
    println math.Floor(3.7)    // 3  (round down)
    println math.Ceil(3.2)     // 4  (round up)
    println math.Round(3.5)    // 4  (nearest)
    println math.Abs(-4.0)     // 4  (absolute value)
end
\end{salam}

\begin{progout}
3
4
4
4
\end{progout}

\subsection{Roots, powers, and logarithms}

\begin{salam}
import "math"

func main:
    println math.Sqrt(144.0)       // 12
    println math.Pow(2.0, 10.0)    // 1024
    println math.Min(2.0, 9.0), math.Max(2.0, 9.0)
end
\end{salam}

\begin{progout}
12
1024
2 9
\end{progout}

\subsection{A reference of common \texttt{math} functions}

\begin{center}
\begin{tabularx}{\linewidth}{@{}lX@{}}
\toprule
\textbf{Function} & \textbf{Returns} \\
\midrule
\code{math.Sqrt(x)} & square root of \code{x} \\
\code{math.Pow(b, e)} & \code{b} raised to the power \code{e} \\
\code{math.Abs(x)} & absolute value (float) \\
\code{math.Floor(x)} / \code{math.Ceil(x)} & round down / up \\
\code{math.Round(x)} & round to nearest \\
\code{math.Min(a, b)} / \code{math.Max(a, b)} & smaller / larger (float) \\
\code{math.MinI(a, b)} / \code{math.MaxI(a, b)} & smaller / larger (integer) \\
\code{math.Gcd(a, b)} / \code{math.Lcm(a, b)} & greatest common divisor / least common multiple \\
\code{math.Sin(x)}, \code{math.Cos(x)}, \code{math.Tan(x)} & trigonometric functions \\
\code{math.Log(x)}, \code{math.Log10(x)}, \code{math.Exp(x)} & logarithms and the exponential \\
\bottomrule
\end{tabularx}
\end{center}

The module also defines the constants \code{math.PI}, \code{math.E}, and
\code{math.TAU}:

\begin{salam}
import "math"

func main:
    r f64 = 2.0
    area auto = math.PI * r * r
    println "area of a circle of radius 2:", area
end
\end{salam}

\begin{progout}
area of a circle of radius 2: 12.5664
\end{progout}

\subsection{Integer helpers}

A handful of functions work directly on integers, so you do not have to cast back
and forth. \code{math.MaxI} and \code{math.MinI} compare integers, and
\code{math.Gcd} and \code{math.Lcm} are exact integer operations:

\begin{salam}
import "math"

func main:
    println math.MaxI(7, 3)    // 7
    println math.Gcd(12, 18)   // 6
    println math.Lcm(4, 6)     // 12
end
\end{salam}

\begin{progout}
7
6
12
\end{progout}

\section{Randomness: \texttt{rand}}

The \code{rand} module produces random values for games, simulations, sampling,
and test data. Start by seeding the generator (use \code{rand.SeedAuto()} for an
unpredictable seed from the clock, or a fixed seed for repeatable runs), then draw
values:

\begin{salam}
import "rand"

func main:
    rand.Seed(42 as i64)                       // fixed seed: repeatable
    dice auto = rand.IntRange(1 as i64, 6 as i64)
    println "you rolled:", dice
    println "coin flip heads?", rand.Bool()
    println "random fraction:", rand.Float()   // in [0, 1)
end
\end{salam}

The most useful functions:

\begin{center}
\begin{tabular}{@{}ll@{}}
\toprule
\textbf{Function} & \textbf{Returns} \\
\midrule
\code{rand.Seed(s)} / \code{rand.SeedAuto()} & set a fixed / clock-based seed \\
\code{rand.IntRange(lo, hi)} & a random integer in \code{[lo, hi]} \\
\code{rand.Float()} & a random \code{f64} in \code{[0, 1)} \\
\code{rand.Bool()} & a random \code{true}/\code{false} \\
\code{rand.UUID()} & a random unique identifier string \\
\bottomrule
\end{tabular}
\end{center}

\begin{tip}
Seed with a \emph{fixed} number while you are developing and testing, so your
program behaves the same on every run and bugs are reproducible. Switch to
\code{rand.SeedAuto()} for the real, shipped program, where you want genuine
unpredictability. Forgetting to seed at all gives you the same ``random'' sequence
every time occasionally useful, usually a surprise.
\end{tip}

\section{Time: \texttt{time}}

The \code{time} module reads the clock, formats timestamps, and pauses your
program. A timestamp is a count of seconds, which you can format into a readable
date or pick apart into fields:

\begin{salam}
import "time"

func main:
    now auto = time.Now()                 // seconds since the epoch
    println "formatted:", time.Format(now)
    println "year:", time.Year(now)
    println "seconds in a day:", time.SecondsPerDay
end
\end{salam}

\code{time.Now()} returns the current time as an integer count of seconds;
\code{time.Format} renders it as \code{"YYYY-MM-DD HH:MM:SS"}, and helpers like
\code{time.Year}, \code{time.Month}, and \code{time.Day} extract individual
fields. The module also defines the constants \code{time.SecondsPerMinute},
\code{time.SecondsPerHour}, and \code{time.SecondsPerDay}.

To pause, use \code{time.Sleep}, which takes a number of milliseconds:

\begin{salam}
import "time"

func main:
    println "wait for it..."
    time.Sleep(500)        // pause half a second
    println "done"
end
\end{salam}

\section{Strings: a reminder}

The \code{str} module \code{Upper}, \code{Lower}, \code{Split}, \code{Join},
\code{Find}, \code{Replace}, \code{FromInt}, \code{ToInt}, and the rest is the
other workhorse of everyday programming, and we covered it thoroughly in
Chapter~8. Keep it in mind alongside the modules here; together, \code{io},
\code{fmt}, \code{math}, \code{str}, \code{rand}, and \code{time} handle the vast
majority of small-program needs.

\section{A worked example: a guessing game}

Here is a complete little program that ties several of these modules together a
number-guessing game (shown with a fixed secret so the output is predictable):

\begin{salam}
import (
    "io"
    "str"
    "rand"
)

func main:
    rand.Seed(7 as i64)
    secret auto = rand.IntRange(1 as i64, 100 as i64)
    io.Print("Guess my number (1-100): ")
    guess auto = str.ToInt(io.Input())
    if guess == secret:
        io.Println("Exactly right!")
    else guess < secret:
        io.Println("Too low.")
    else:
        io.Println("Too high.")
    end
end
\end{salam}

It seeds the generator, picks a secret in $[1, 100]$, reads a guess as text,
converts it to a number, and compares drawing on \code{rand}, \code{io}, and
\code{str} at once. That is the standard library doing exactly what it is for:
letting you assemble a real program from ready-made parts.

\section{Chapter summary}

\begin{itemize}[leftmargin=1.4em]
  \item \code{io} provides \code{Print}, \code{Println}, and \code{Input} (read a
        line as a \code{str}).
  \item \code{fmt} converts values (\code{Int}, \code{Float}, \code{Bool}) and
        aligns text (\code{PadLeft}, \code{PadRight}, \code{Repeat}).
  \item \code{math} offers rounding, roots, powers, trig, logs, integer helpers
        (\code{MaxI}, \code{Gcd}, \code{Lcm}), and the constants \code{PI},
        \code{E}, \code{TAU}.
  \item \code{rand} generates random integers, floats, booleans, and ids; seed it
        for repeatability or with \code{SeedAuto} for unpredictability.
  \item \code{time} reads and formats the clock, extracts date fields, exposes
        second-count constants, and sleeps with \code{Sleep}.
\end{itemize}

\section{Exercises}

\exercise Write a program that asks for the user's name and age and prints a
sentence using both (convert the age with \code{str.ToInt}).

\exercise Print a small multiplication table aligned in columns using
\code{fmt.PadLeft}.

\exercise Use \code{math} to print the hypotenuse of a right triangle with legs 3
and 4 (\code{Sqrt} of the sum of squares).

\exercise Roll two dice with \code{rand.IntRange} and print their sum; seed with a
fixed value so you can predict the result.

\exercise Print the current date with \code{time.Format(time.Now())}, then sleep
one second and print it again.

\exercise Combine \code{rand} and \code{fmt} to print five random integers between
1 and 99, each padded to width 2, on one line.
